% META: {"id":"TCOMPL-ME-ME-004","chapitre":"TCOMPL-MODELES-EVOLUTION","type_objet":"methode","capacites_codes":["C4"],"methodes":["M4"],"mode_creation":"ex_nihilo","sources_inspiration":[],"status":"approved","origin":"needs_review","status_history":["needs_review","audited","accepted","approved"],"release_acceptance":"RELEASE_OWNER_BATCH_ACCEPTANCE","acceptance_closure_digest":"sha256:9b3ccf9a81c5520fbb7e03b7b2d3e7bf2057a02834b6d908bf3be5f49f3b553f","acceptance_receipt_digest":"sha256:4fad86f0282e0fa4e0419cca1ad6379ae7af5a280c04a4a90880f90a1c044242"}
% BEGIN-VERIFY
% from sympy import symbols, Rational, limit, oo, simplify
% n = symbols('n', integer=True, nonnegative=True)
% un = 4 - 3*Rational(1,2)**n
% assert un.subs(n, 0) == 1
% assert simplify(un.subs(n, n+1) - (un/2 + 2)) == 0
% assert limit(un, n, oo) == 4
% END-VERIFY
\begin{fichemethode}{M4}{Résoudre une suite arithmetico-géométrique via sa solution constante}

\quandUtiliser{%
  Utiliser cette méthode lorsque la suite vérifie $u_{n+1} = a\,u_n + b$
  (avec $a \neq 1$) et que l'énoncé demande une formule explicite de $u_n$
  et/ou sa limite.
}

\pasApas{%
  \begin{enumerate}
    \item Chercher la solution CONSTANTE : résoudre $c = a\,c + b$, soit
      $c = \dfrac{b}{1-a}$.
    \item Poser la suite auxiliaire $v_n = u_n - c$.
    \item Montrer que $(v_n)$ est géométrique :
      $v_{n+1} = u_{n+1} - c = a\,u_n + b - (ac + b) = a(u_n - c) = a\,v_n$.
    \item En déduire $v_n = v_0\,a^n = (u_0 - c)\,a^n$, puis la formule
      explicite $u_n = c + (u_0 - c)\,a^n$.
    \item Conclure sur la limite : si $|a| < 1$, $a^n \to 0$ donc $u_n \to c$ —
      la solution constante est la valeur d'equilibre.
  \end{enumerate}
}

\exempleRedige{%
  Soit $u_0 = 1$ et $u_{n+1} = \dfrac{1}{2}\,u_n + 2$. Exprimer $u_n$ en
  fonction de $n$ et déterminer sa limite.

  \medskip
  \commentaireMarge{Point fixe : $c = ac + b$.}%
  Solution constante : $c = \dfrac{c}{2} + 2$ donne $c = 4$.

  \medskip
  \commentaireMarge{Suite auxiliaire centree sur l'equilibre.}%
  Posons $v_n = u_n - 4$. Alors
  $v_{n+1} = u_{n+1} - 4 = \dfrac{u_n}{2} + 2 - 4 = \dfrac{u_n - 4}{2}
  = \dfrac{v_n}{2}$ : $(v_n)$ est géométrique de raison $\dfrac{1}{2}$ et de
  premier terme $v_0 = 1 - 4 = -3$.

  \medskip
  Donc $v_n = -3\left(\dfrac{1}{2}\right)^n$ et
  $u_n = 4 - 3\left(\dfrac{1}{2}\right)^n$. Comme
  $\left(\frac{1}{2}\right)^n \to 0$, $\lim u_n = 4$.
}

\pieges{%
  \begin{itemize}
    \item \textbf{Piège 1.} Se tromper de suite auxiliaire : c'est $u_n - c$
      (ecart a l'equilibre), pas $u_n - u_0$.
    \item \textbf{Piège 2.} Oublier de calculer $v_0$ avec $u_0$ : la formule
      finale depend de $v_0 = u_0 - c$.
    \item \textbf{Piège 3.} Conclure $u_n \to c$ sans avoir vérifie $|a| < 1$
      (si $|a| > 1$ et $u_0 \neq c$, la suite diverge).
  \end{itemize}
}

\verifier{%
  Contrôler la formule explicite sur $u_0$, $u_1$, $u_2$ (récurrence directe),
  vérifier qu'elle satisfait la relation $u_{n+1} = a\,u_n + b$, et confronter
  la limite a la conjecture graphique (fiche M3).
}

\sEntrainer{\refExos{M4}}

\end{fichemethode}
